% gongbu-haja DEEP 노트 그림 예시 — notefigure 환경과 그림 스타일을 쓰는 대표 유형 다섯 가지.
%
% 노트에 그림을 처음 넣거나 새 유형(그래프·힘 도식·블록선도 등)을 그릴 때 해당 예시만 열어 본뜬다.
% 디자인 값(색·선 굵기·글꼴·캡션)은 deep_note_style.tex가 정하므로 여기 값을 노트 서문에 옮겨 적지 않는다.
% 이 파일은 테스트에서 실제로 조판된다: scripts 폴더에서 xelatex -no-shell-escape deep_figure_gallery.tex
%
%   1. 함수 그래프와 핵심 점      pgfplots axis, 축 끝 라벨, 핵심 값 눈금, 보조선, 점
%   2. 곡선 두 개 비교            순서대로 남색 실선·검정 파선, 곡선 끝 이름, 꼭짓점 간격 표시(dimension)
%   3. 넓이(적분 구간)            구간만 칠하기, 구간 끝 표시, 영역 안 수식
%   4. 힘 도식(빗면)              tikz 좌표 회전, vector, 성분 분해, 각도 표시
%   5. 블록선도(되먹임)           block·sum·signal, positioning
\documentclass[a4paper,11pt]{article}
\input{"deep_note_style.tex"}
\gongbuheader{그림 예시}{노트가 그리는 그림}
\begin{document}

\section{함수 그래프와 핵심 점}

축전기가 충전될 때 전압은 $v_C(t)=V_0\bigl(1-e^{-t/\tau}\bigr)$를 따른다. 그림~\ref{fig:rc}에서 볼 것은 두 가지다.
곡선은 처음에 가장 가파르게 오르다가 점점 완만해지며 $V_0$에 다가가기만 하고 닿지는 않는다. 그리고
$t=\tau$에서 이미 최종값의 $1-e^{-1}\approx 0.632$, 곧 약 63\,\%에 도달한다.

% 축 범위·눈금은 본문이 말하는 값(τ, 0.632V0, V0)에 맞춘다. 눈금 숫자는 수식 글꼴로 조판된다.
% 곡선은 \addplot+[…]로 그려 정해진 색·선 모양 순서(남색 실선 → 검정 파선 → 갈색 일점쇄선)를 따른다.
% \addplot[…]처럼 +를 빼면 그 순서가 꺼지므로 보조선(guide)·영역(region)에만 쓴다.
\begin{notefigure}{축전기 충전 전압. $t=\tau$에서 최종값의 약 63\,\%에 이른다.}
\begin{tikzpicture}
\begin{axis}[xmin=0,xmax=5,ymin=0,ymax=1.1,
  xlabel={$t$},ylabel={$v_C(t)$},
  xtick={1,2,3,4,5},xticklabels={$\tau$,$2\tau$,$3\tau$,$4\tau$,$5\tau$},
  ytick={0.632,1},yticklabels={$0.632V_0$,$V_0$}]
  \addplot[guide,domain=0:5] {1};
  \addplot+[domain=0:5] {1-exp(-x)};
  \draw[guide] (axis cs:1,0) |- (axis cs:0,0.632);
  \node[dot=FigMain] at (axis cs:1,0.632) {};
\end{axis}
\end{tikzpicture}
\label{fig:rc}
\end{notefigure}

\section{곡선 두 개 비교}

교류 회로에서 전압 $v(t)=V_m\sin\omega t$보다 전류 $i(t)=I_m\sin(\omega t-\varphi)$가 늦으면, 그림~\ref{fig:phase}처럼
전류의 꼭짓점이 전압의 꼭짓점보다 오른쪽에 온다. 두 꼭짓점 사이의 가로 거리가 위상차 $\varphi$를 각도로 나타낸 것이며,
두 곡선의 높이 차이는 위상차와 관계없는 진폭 $V_m$과 $I_m$의 차이다.

% 곡선이 둘 이상이면 범례 대신 곡선 끝에 이름을 단다(node[pos=…]). 이름끼리 겹칠 때만 범례를 쓴다.
% 가운데 축을 곡선이 자주 가로지르면 그 축의 눈금 숫자가 곡선과 겹치므로 필요 없는 눈금은 뺀다.
\begin{notefigure}{전류가 전압보다 $\varphi$만큼 늦은 경우. 꼭짓점 사이의 거리가 위상차다.}
\begin{tikzpicture}
\begin{axis}[xmin=0,xmax=7,ymin=-1.15,ymax=1.35,
  xlabel={$\omega t$},xtick=\empty,ytick=\empty]
  \addplot+[domain=0:2*pi] {sin(deg(x))} node[pos=0.12,anchor=south east] {$v(t)$};
  \addplot+[domain=0:2*pi] {0.7*sin(deg(x-pi/3))} node[pos=0.55,anchor=south west] {$i(t)$};
  \draw[guide] (axis cs:pi/2,1) -- (axis cs:pi/2,1.22);
  \draw[guide] (axis cs:5*pi/6,0.7) -- (axis cs:5*pi/6,1.22);
  \draw[dimension] (axis cs:pi/2,1.17) -- node[above] {$\varphi$} (axis cs:5*pi/6,1.17);
\end{axis}
\end{tikzpicture}
\label{fig:phase}
\end{notefigure}

\section{넓이(적분 구간)}

정적분 $\int_a^b f(x)\,dx$는 그림~\ref{fig:area}에서 곡선 아래, $x=a$와 $x=b$ 사이의 칠한 부분의 넓이다.
구간 밖의 곡선은 넓이에 들어가지 않는다는 점을 칠한 영역의 경계로 바로 확인할 수 있다.

% 구간만 칠할 때는 같은 식을 domain=a:b로 한 번 더 그리고 \closedcycle로 x축까지 닫는다.
\begin{notefigure}{곡선 아래 $a\le x\le b$ 구간의 넓이가 정적분이다.}
\begin{tikzpicture}
\begin{axis}[xmin=0,xmax=3.2,ymin=0,ymax=3.2,
  xlabel={$x$},ylabel={$y$},xtick={1,2.5},xticklabels={$a$,$b$},ytick=\empty]
  \addplot[region,domain=1:2.5] {0.25*x^2+0.6} \closedcycle;
  \addplot+[domain=0:3.1] {0.25*x^2+0.6} node[pos=0.97,anchor=south east] {$y=f(x)$};
  \draw[guide] (axis cs:1,0) -- (axis cs:1,0.85);
  \draw[guide] (axis cs:2.5,0) -- (axis cs:2.5,2.1625);
  \node at (axis cs:1.75,0.6) {$\displaystyle\int_a^b f(x)\,dx$};
\end{axis}
\end{tikzpicture}
\label{fig:area}
\end{notefigure}

\section{힘 도식}

빗면 위 물체에 작용하는 중력 $mg$를 빗면 방향과 빗면에 수직인 방향으로 나누면 그림~\ref{fig:incline}의
두 파선 성분이 된다. 빗면을 따라 내려가는 성분은 $mg\sin\theta$, 빗면을 누르는 성분은 $mg\cos\theta$이며,
수직항력 $N$은 이 누르는 성분과 크기가 같고 방향이 반대다.

% 빗면 방향의 그림은 scope의 rotate로 좌표계를 돌려 그린다. 길이 비율은 실제 식(sin·cos)에 맞춘다.
\begin{notefigure}{빗면 위 물체의 힘. 중력을 빗면 방향과 수직 방향 성분으로 나눈다.}
\begin{tikzpicture}[scale=1.05]
  \coordinate (O) at (0,0);
  \coordinate (B) at (6,0);
  \coordinate (T) at (6,3.464);
  \draw[body,fill=Hairline!45] (O) -- (B) -- (T) -- cycle;
  \pic[draw=FigInk,line width=0.45pt,angle radius=9mm,angle eccentricity=1.3,"$\theta$"] {angle=B--O--T};
  \begin{scope}[shift={(3.464,2)},rotate=30]
    \draw[body] (-0.55,0) rectangle (0.55,0.75);
    \coordinate (C) at (0,0.375);
    \draw[vector=FigMain,densely dashed] (C) -- ++(-0.8,0) node[above left=-1pt] {$mg\sin\theta$};
    \draw[vector=FigMain,densely dashed] (C) -- ++(0,-1.386) node[below right=-1pt] {$mg\cos\theta$};
    \draw[vector] (C) -- ++(0,1.386) node[above] {$N$};
  \end{scope}
  \draw[vector] (C) -- ++(0,-1.6) node[left=1pt,pos=0.9] {$mg$};
\end{tikzpicture}
\label{fig:incline}
\end{notefigure}

\section{블록선도}

되먹임 제어에서 출력 $C(s)$는 $H(s)$를 거쳐 입력 쪽으로 돌아와 기준 $R(s)$에서 빼진다. 그림~\ref{fig:loop}의
합산점에서 이 차이 $E(s)=R(s)-H(s)C(s)$가 만들어지고, $C(s)=G(s)E(s)$이므로 전달함수는
$C(s)/R(s)=G(s)/\bigl(1+G(s)H(s)\bigr)$이다.

% 노드 사이 간격은 positioning(right=… of …)으로 정하고, 합산점의 부호는 들어오는 선 옆에 단다.
\begin{notefigure}{음의 되먹임 블록선도. 합산점에서 되먹임 신호를 뺀다.}
\begin{tikzpicture}
  \node (in) {$R(s)$};
  \node[sum,right=10mm of in] (sum) {};
  \node[block,right=12mm of sum] (G) {제어 대상\\$G(s)$};
  \coordinate[right=10mm of G] (split);
  \node[right=8mm of split] (out) {$C(s)$};
  \node[block,below=6mm of G] (H) {센서 $H(s)$};
  \draw[signal] (in) -- node[pos=0.85,above] {$+$} (sum);
  \draw[signal] (sum) -- node[above] {$E(s)$} (G);
  \draw[signal] (G) -- (out);
  \draw[signal] (split) |- (H);
  \draw[signal] (H) -| node[pos=0.92,left] {$-$} (sum);
  \node[dot,inner sep=0pt,minimum size=4pt] at (split) {};
\end{tikzpicture}
\label{fig:loop}
\end{notefigure}

\end{document}
